% !TEX spellcheck = en_US

% ~~~~~~~~~~~~~~~~~~~~~~~~~~~~~~~~~~~~~~~ V
% (NC) 2026 Haluk Bingol
% github.com/halukbingol/LaTeX-Templates
%
% Licensees may copy, distribute, display, and perform the work and make derivative works
% and remixes based on it only for non-commercial purposes.
% ~~~~~~~~~~~~~~~~~~~~~~~~~~~~~~~~~~~~~~~ A

% v20261108T0905   the libraries are split: articles load hblibArticle (which loads hblibCommon);
%                  \usepackage{hblibCommon} still works; "Getting started" lists the five libraries
% v20261107T1235   hyperlinks: boxes (hblibCommon v20261107T1230)
% v20261106T1210   \title[short]{...}: TeX Live 2026 looped ("TeX capacity exceeded") on \normalsize in the
%                  page-header copy of the title; troubleshooting row added
% v20261106T1030   user manual: quick reference (all options, commands, styles, colors),
%                  customizing, troubleshooting; title "hblibCommon User Manual"
% v20261106T1010   renamed hblibCommon-test.tex to hblibCommon-UserManual.tex
% v20261106T0945   lstlisting blocks wrapped in V / A markers too (with their \hCodeTag line)
% v20261106T0905   code blocks (texcode, javacode, bashcode) wrapped in V / A markers
% v20261106T0830   examples and notes for hblibCommon.sty v20261107T1230
%
% hblibCommon-UserManual.tex -- the user manual of hblibCommon.sty: a quick reference, then every feature
% with the LaTeX source next to the result, customizing and troubleshooting.
% Put it next to hblibArticle.sty and hblibCommon.sty (or change the \usepackage line) and compile twice:
%   pdflatex hblibCommon-UserManual.tex
%   pdflatex hblibCommon-UserManual.tex
% It writes two small example files, D.java and run.sh, next to itself (see "Example files" below).

\documentclass[11pt]{amsart}         % article works too



\usepackage{hblibArticle}            % from another folder: \usepackage{../../cis000_latex/lib/hblibArticle}



% +++++++++++++++++++++++++++++++++++++++ V specific to this test
% texcode: shows LaTeX source in a box tagged "latex" (it can contain \end{lstlisting}, \end{javacode}, ...)
\lstnewenvironment{texcode}{\hCodeTag{latex}\lstset{style=texsnippet, numbers=none}}{}
% \hNote{text}: a note in a light box
\newcommand{\hNote}[1]{\par\smallskip\noindent\colorbox{codebg}{\parbox{\dimexpr\linewidth-2\fboxsep}{\small\raggedright\textbf{Note.} #1}}\par\smallskip}
% \cs{name}: a command name \name;  \marg{x}: {x};  \oarg{x}: [x]   (in the reference tables)
\newcommand{\cs}[1]{\hCode{\textbackslash #1}}
\newcommand{\marg}[1]{\hCode{\{#1\}}}
\newcommand{\oarg}[1]{\hCode{[#1]}}
% +++++++++++++++++++++++++++++++++++++++ A specific to this test



% ~~~~~~~~~~~~~~~~~~~~~~~~~~~~~~~~~~~~~~~ V example files
% filecontents* writes the files used by \javafile and \bashfile below ([overwrite]: every run)
\begin{filecontents*}[overwrite]{D.java}
package demo;

/** D shows how line numbers work with firstline. */
public class D {
    public static void main(String[] args) {
        int n = args.length;          // line 6
        System.out.println("n = " + n);
    }
}
\end{filecontents*}

\begin{filecontents*}[overwrite]{run.sh}
#!/bin/bash
# compile and run D
javac -d out D.java
java -cp out demo.D aaa bbb
\end{filecontents*}
% ~~~~~~~~~~~~~~~~~~~~~~~~~~~~~~~~~~~~~~~ A example files



% ~~~~~~~~~~~~~~~~~~~~~~~~~~~~~~~~~~~~~~~ V title
\title[hblibCommon User Manual]{%   [short title]: the page-header copy must not contain size commands
	hblibCommon User Manual\\[2pt]
	{\normalsize for hblibArticle.sty and hblibCommon.sty v20261108T0900}
}
\author{Haluk Bingol}

\begin{document}

\maketitle
\tableofcontents
% ~~~~~~~~~~~~~~~~~~~~~~~~~~~~~~~~~~~~~~~ A title




% ~~~~~~~~~~~~~~~~~~~~~~~~~~~~~~~~~~~~~~
\section{Getting started}\label{sec:loading}

% +++++++++++++++++++++++++++++++++++++++ V
%: -Code.
\begin{texcode}
\documentclass[11pt]{amsart}           % or article; [twocolumn] works too
\usepackage{hblibArticle}              % titles and headings as typed
\usepackage[amsheadings]{hblibArticle} % amsart: keep CAPITAL titles, small-caps headings
\usepackage{../../cis000_latex/lib/hblibArticle}   % from another folder
\usepackage{../../cis000_latex/lib/hblibDatastructures} % optional: arrays, stacks, trees
\usepackage{../../cis000_latex/lib/hblibCodeoutput}     % optional: code and output from files
\end{texcode}
% +++++++++++++++++++++++++++++++++++++++ A

\begin{center}\small
\begin{tabular}{@{}L{0.27\linewidth}L{0.67\linewidth}@{}}
	\toprule
	\textbf{Library} & \textbf{Contents} \\
	\midrule
	hblibCommon & shared by articles and slides: fonts, packages, algorithms, colors and emphasis,
		code styles and environments, \cs{hCode}, time stamps, \cs{reffig} \dots, \cs{hAbs}, column type \hCode{L};
		not loaded by the document, but by the two below \\
	hblibArticle & articles: page, lists, titles as typed, contents, section numbering, footer, hyperlinks \\
	hblibPresentation & slides (beamer): theme, \cs{hSection}, the old \hCode{mystyle} listings \\
	hblibDatastructures & optional: arrays, stacks, queues, strings, grids, trees, commit graphs \\
	hblibCodeoutput & optional: code and output read from files, in boxes (tcolorbox) \\
	\bottomrule
\end{tabular}
\end{center}

\begin{itemize}\raggedright
	\item Keep the five files in one folder: each library loads the others it needs from its own folder.
	\item Old documents with \hCode{\textbackslash usepackage\{hblibCommon\}} need no change:
	hblibCommon then loads hblibArticle (the log says so).
	\item Load it with \hCode{\textbackslash usepackage}, without \hCode{.sty}.
	\hCode{\textbackslash input\{hblibCommon.sty\}} does \emph{not} work:
	names with \hCode{@} inside the library break, and LaTeX reports ``Missing \textbackslash begin\{document\}''.
	\item With a relative path LaTeX warns
	``You have requested package `../../cis000\_latex/hblibCommon', but the package provides `hblibCommon'''.
	The warning is harmless.
	\item Compile \emph{twice}: the contents, the references and the ``page / pages'' in the footer
	are known only after the first run.
	\item Already loaded, so do not load again:
	fontenc (T1), lmodern, geometry, amsmath, amssymb, graphicx, xcolor, tikz, calc, array, booktabs, enumitem,
	forest, datetime2, menukeys, algorithm2e, listings, fancyhdr, lastpage, tocloft (not with amsart), hyperref.
	\item Page: A4 with 2\,cm margins. Change it after \hCode{\textbackslash usepackage\{hblibArticle\}}, e.g.\
	\hCode{\textbackslash geometry\{bottom=2.5cm\}}.
\end{itemize}




% ~~~~~~~~~~~~~~~~~~~~~~~~~~~~~~~~~~~~~~
\section{Quick reference}\label{sec:quickref}

Every option, command and environment of hblibCommon, with the section that shows an example.
Optional arguments are in \oarg{brackets}.

% ~~~~~~~~~~~~~~~~~~~~~~~~~~~~~~~~~~~~~~
\subsection{Option}

\begin{center}\small
\begin{tabular}{@{}L{0.36\linewidth}L{0.58\linewidth}@{}}
	\toprule
	\textbf{Write} & \textbf{Result} \\
	\midrule
	\cs{usepackage}\marg{hblibArticle} & titles, authors and headings as typed (also with amsart) \\
	\cs{usepackage}\oarg{amsheadings}\marg{hblibArticle} & amsart: keep its CAPITAL title and small-caps headings \\
	\bottomrule
\end{tabular}
\end{center}

% ~~~~~~~~~~~~~~~~~~~~~~~~~~~~~~~~~~~~~~
\subsection{Code}

\begin{center}\small
\begin{tabular}{@{}L{0.36\linewidth}L{0.58\linewidth}@{}}
	\toprule
	\textbf{Write} & \textbf{Result} \\
	\midrule
	\hCode{\textbackslash begin\{javacode\}}\oarg{options} & typed Java: framed box, tag ``java'' \\
	\hCode{\textbackslash begin\{bashcode\}}\oarg{options} & typed bash: framed box, tag ``bash'' \\
	\cs{javafile}\oarg{options}\marg{file} & Java from a file, line numbers, tag ``java'' \\
	\cs{bashfile}\oarg{options}\marg{file} & bash from a file, tag ``bash'' \\
	\cs{hCode}\marg{text} & inline code: typewriter, colored, 2\,pt smaller; escape \hCode{\_ \{ \} \textbackslash} \\
	\cs{hCodeJ}\marg{code} & inline Java, syntax colored, typed as is \\
	\cs{hCodeB}\marg{code} & inline bash, typed as is \\
	\cs{code}\marg{code} & old name of \cs{hCodeJ} \\
	\cs{hCodeTag}\marg{text} & the next listing shows \emph{text} in its upper right corner \\
	\cs{hCodeTagsfalse}, \cs{hCodeTagstrue} & no tags from here on; tags again \\
	\hCode{style=}\emph{name} & listing styles: \reftbl{tbl:styles} \\
	\bottomrule
\end{tabular}
\end{center}

{\raggedright
Options are ordinary listings options, e.g.\ \hCode{numbers=left}, \hCode{firstline=5}, \hCode{lastline=7},
\hCode{firstnumber=5}, \hCode{basicstyle=\textbackslash ttfamily \textbackslash scriptsize}.
Examples: \refsec{sec:code}.\par}

% ~~~~~~~~~~~~~~~~~~~~~~~~~~~~~~~~~~~~~~
\subsection{Section numbering}

\begin{center}\small
\begin{tabular}{@{}L{0.36\linewidth}L{0.58\linewidth}@{}}
	\toprule
	\textbf{Write} & \textbf{Result} \\
	\midrule
	\cs{hNumbersLetter}\marg{G} & next section G; subsections G1, G2; subsubsections G1.1 \\
	\cs{hNumbersArabic} & back to 1, 1.1; the numbers continue where they stopped \\
	\cs{hNumbersArabic}\oarg{1} & back to numbers; the next section is 1 \\
	\cs{hSectionWarmUp} & section ``W WarmUp'' \\
	\cs{hSectionGoal} & section ``G Goal'' \\
	\cs{hSectionChallenge} & section ``C Challenge'' \\
	\cs{hSectionHome} & section ``H Home'' \\
	\bottomrule
\end{tabular}
\end{center}

Example: \refsec{sec:numbering}.

% ~~~~~~~~~~~~~~~~~~~~~~~~~~~~~~~~~~~~~~
\subsection{References}

\begin{center}\small
\begin{tabular}{@{}L{0.30\linewidth}L{0.20\linewidth}L{0.30\linewidth}L{0.14\linewidth}@{}}
	\toprule
	\textbf{Write} & \textbf{Result} & \textbf{Write} & \textbf{Result} \\
	\midrule
	\cs{reffig}\marg{label} & Fig.~3        & \cs{refthm}\marg{label} & Theorem~1 \\
	\cs{refeq}\marg{label}  & Eq.~2         & \cs{refthmA}\marg{thm}\marg{item} & Theorem~1(b) \\
	\cs{reftbl}\marg{label} & Table~1       & \cs{reflem}\marg{label} & Lemma~2 \\
	\cs{refsec}\marg{label} & Sec.~4        & \cs{refdef}\marg{label} & Definition~3 \\
	\cs{refalg}\marg{label} & Alg.~1        & \cs{refexmp}\marg{label} & Example~4 \\
	\cs{reflst}\marg{label} & List~2        & \cs{refcite}\marg{key} & Ref.~[7] \\
	\bottomrule
\end{tabular}
\end{center}

% ~~~~~~~~~~~~~~~~~~~~~~~~~~~~~~~~~~~~~~
\subsection{Text and time stamps}

\begin{center}\small
\begin{tabular}{@{}L{0.36\linewidth}L{0.58\linewidth}@{}}
	\toprule
	\textbf{Write} & \textbf{Result} \\
	\midrule
	\cs{hQuote}\marg{text} & a short quotation: small, sans serif, in quotes \\
	\cs{hIdea}\marg{text} & a note to self: tiny, olive, in brackets \\
	\cs{hFootnote}\marg{who}\marg{text} & a footnote starting with a red ``@who :'' \\
	\cs{hAbs}\marg{x} & $\lvert x \rvert$ with bars that fit \\
	\cs{hFontLarger}\marg{N}\marg{text} & text $N$\,pt larger ($N<0$: smaller); old name \cs{hbFontShift} \\
	\cs{YULearn} & \YULearn, the course web site \\
	\cs{hbTimeStamp} & v\today, used in the footer \\
	\cs{hbTimeStampDraft} & red ``--Draft--'' with date and time \\
	\bottomrule
\end{tabular}
\end{center}

% ~~~~~~~~~~~~~~~~~~~~~~~~~~~~~~~~~~~~~~
\subsection{Algorithms}

algorithm2e is loaded with \hCode{ruled, vlined, linesnumbered} and captions in \hCode{\textbackslash small}.
hblibCommon adds \cs{Fn}\marg{name}\marg{body} (``function \dots\ end'') and the aligned
\cs{Input}\marg{text} and \cs{Output}\marg{text} lines.
Example: \refalg{alg:balanced}.

% ~~~~~~~~~~~~~~~~~~~~~~~~~~~~~~~~~~~~~~
\subsection{Colors}\label{sec:colors}

All can be changed with \cs{definecolor}, see \refsec{sec:customizing}.

\begin{center}\small
\begin{tabular}{@{}L{0.20\linewidth}L{0.20\linewidth}L{0.52\linewidth}@{}}
	\toprule
	\textbf{Name} & \textbf{RGB} & \textbf{Used for} \\
	\midrule
	colorHCode & 240, 80, 50   & \cs{hCode} \\
	codebg     & 248, 248, 248 & background of code boxes \\
	codekw     & 0, 70, 160    & keywords \\
	codestr    & 30, 130, 60   & strings \\
	codecmt    & 120, 120, 120 & comments, line numbers, tags \\
	codeframe  & 40, 90, 160   & code box frame (as in the slide decks) \\
	termbg, termfg, termframe & 35, 38, 45 \dots & terminal: background, text, frame \\
	termcmd    & 255, 214, 102 & terminal: lines starting with \$ \\
	\bottomrule
\end{tabular}
\end{center}




% ~~~~~~~~~~~~~~~~~~~~~~~~~~~~~~~~~~~~~~
\section{Code}\label{sec:code}

% ~~~~~~~~~~~~~~~~~~~~~~~~~~~~~~~~~~~~~~
\subsection{Typed code: javacode, bashcode}

% +++++++++++++++++++++++++++++++++++++++ V
%: -Code.
\begin{texcode}
\begin{javacode}
public class Hello {   // Turkish letters work: ğüşıöç
    public static void main(String[] args) {
        System.out.println("Hello");
    }
}
\end{javacode}
\end{texcode}
% +++++++++++++++++++++++++++++++++++++++ A

% +++++++++++++++++++++++++++++++++++++++ V
%: -Code.
\begin{javacode}
public class Hello {   // Turkish letters work: ğüşıöç
    public static void main(String[] args) {
        System.out.println("Hello");
    }
}
\end{javacode}
% +++++++++++++++++++++++++++++++++++++++ A

% +++++++++++++++++++++++++++++++++++++++ V
%: -Code.
\begin{texcode}
\begin{bashcode}
javac Hello.java   # compile
java Hello         # run
\end{bashcode}
\end{texcode}
% +++++++++++++++++++++++++++++++++++++++ A

% +++++++++++++++++++++++++++++++++++++++ V
%: -Code.
\begin{bashcode}
javac Hello.java   # compile
java Hello         # run
\end{bashcode}
% +++++++++++++++++++++++++++++++++++++++ A

Any listings option can be given in brackets, e.g.\ line numbers:
\hCode{\textbackslash begin\{javacode\}[numbers=left]}.

% +++++++++++++++++++++++++++++++++++++++ V
%: -Code.
\begin{javacode}[numbers=left]
int a = 1;
int b = a + 1;
\end{javacode}
% +++++++++++++++++++++++++++++++++++++++ A


% ~~~~~~~~~~~~~~~~~~~~~~~~~~~~~~~~~~~~~~
\subsection{Code from a file: \textbackslash javafile, \textbackslash bashfile}

The whole file, with line numbers (\hCode{\textbackslash javafile} numbers the lines by default):

% +++++++++++++++++++++++++++++++++++++++ V
%: -Code.
\begin{texcode}
\javafile{D.java}
\end{texcode}
% +++++++++++++++++++++++++++++++++++++++ A

\javafile{D.java}

% +++++++++++++++++++++++++++++++++++++++ V
%: -Code.
\begin{texcode}
\bashfile{run.sh}
\end{texcode}
% +++++++++++++++++++++++++++++++++++++++ A

\bashfile{run.sh}


% ~~~~~~~~~~~~~~~~~~~~~~~~~~~~~~~~~~~~~~
\subsection{Part of a file and its line numbers}\label{sec:firstnumber}

\hNote{\hCode{\textbackslash javafile[firstline=5, lastline=7]\{D.java\}}
numbers the lines 1, 2, 3, not 5, 6, 7.
Add \hCode{firstnumber=5} to the options when you want the file's own line numbers.}

% +++++++++++++++++++++++++++++++++++++++ V
%: -Code.
\begin{texcode}
\javafile[firstline=5, lastline=7]{D.java}                 % numbered 1, 2, 3
\javafile[firstline=5, lastline=7, firstnumber=5]{D.java}  % numbered 5, 6, 7
\end{texcode}
% +++++++++++++++++++++++++++++++++++++++ A

\javafile[firstline=5, lastline=7]{D.java}

\javafile[firstline=5, lastline=7, firstnumber=5]{D.java}

Now the text can say ``line 6'' and mean the same line as the file.


% ~~~~~~~~~~~~~~~~~~~~~~~~~~~~~~~~~~~~~~
\subsection{Language tags}\label{sec:tags}

As in markdown, a code box shows its language small and gray in its upper right corner.
\begin{itemize}\raggedright
	\item \hCode{bashcode} and \hCode{\textbackslash bashfile} show \textbf{bash}.
	\item \hCode{javacode} and \hCode{\textbackslash javafile} show \textbf{java},
	also when you print part of a file with \hCode{firstline=\dots}.
	\item Other listings: put \hCode{\textbackslash hCodeTag\{\dots\}} right before the listing.
	For example, \hCode{\textbackslash hCodeTag\{txt\}} before a \hCode{MANIFEST.MF} box
	matches a markdown ``txt'' block.
	\item Turning them off: \hCode{\textbackslash hCodeTagsfalse} in a document turns tags off from that point on;
	\hCode{\textbackslash hCodeTagstrue} turns them on again.
\end{itemize}

% +++++++++++++++++++++++++++++++++++++++ V
%: -Code.
\begin{texcode}
\hCodeTag{txt}
\begin{lstlisting}[style=outputstyle, frame=single, rulecolor=\color{gray!40}]
Main-Class: MyClass
\end{lstlisting}
\end{texcode}
% +++++++++++++++++++++++++++++++++++++++ A

% +++++++++++++++++++++++++++++++++++++++ V
%: -Code.
\hCodeTag{txt}
\begin{lstlisting}[style=outputstyle, frame=single, rulecolor=\color{gray!40}]
Main-Class: MyClass
\end{lstlisting}
% +++++++++++++++++++++++++++++++++++++++ A

% +++++++++++++++++++++++++++++++++++++++ V
%: -Code.
\begin{texcode}
\hCodeTagsfalse
\begin{bashcode}
ls -l    # no tag
\end{bashcode}
\hCodeTagstrue
\end{texcode}
% +++++++++++++++++++++++++++++++++++++++ A

\hCodeTagsfalse
% +++++++++++++++++++++++++++++++++++++++ V
%: -Code.
\begin{bashcode}
ls -l    # no tag
\end{bashcode}
% +++++++++++++++++++++++++++++++++++++++ A
\hCodeTagstrue

\hNote{The tag takes no space of its own: it is drawn at the right end of the first line.
A first line that runs the full width of the box goes under the tag.}


% ~~~~~~~~~~~~~~~~~~~~~~~~~~~~~~~~~~~~~~
\subsection{Other listing styles}

All of them are used as \hCode{\textbackslash begin\{lstlisting\}[style=\dots]}.

% +++++++++++++++++++++++++++++++++++++++ V tbl
%: -Tbl. styles
\begin{table}[htbp]
\caption{Listing styles of hblibCommon.}
\label{tbl:styles}
\small
\begin{tabular}{@{}>{\ttfamily}l >{\raggedright\arraybackslash}p{0.62\linewidth}@{}}
	\toprule
	\textrm{\textbf{Style}} & \textbf{Use} \\
	\midrule
	codesnippet   & short code typed in the .tex: framed, light background (javacode and bashcode use it) \\
	codestyle     & colored code with line numbers, no frame (the default of a plain lstlisting) \\
	outputstyle   & program output: no colors, no numbers \\
	terminalstyle & a terminal session: lines starting with \$ are highlighted \\
	texsnippet    & LaTeX source \\
	javastyle, bashstyle & codesnippet plus the language \\
	mystyle       & old name (hblibSimple); now the same as codestyle \\
	\bottomrule
\end{tabular}
\end{table}%
% +++++++++++++++++++++++++++++++++++++++ A tbl

\reftbl{tbl:styles} lists them. A terminal session:

% +++++++++++++++++++++++++++++++++++++++ V
%: -Code.
\begin{texcode}
\hCodeTag{terminal}
\begin{lstlisting}[style=terminalstyle, frame=single, rulecolor=\color{termframe}, backgroundcolor=\color{termbg}, basicstyle=\ttfamily\footnotesize\color{termfg}]
$ java -jar MyApp.jar aaa
~~MyClass
i=0:aaa
\end{lstlisting}
\end{texcode}
% +++++++++++++++++++++++++++++++++++++++ A

% +++++++++++++++++++++++++++++++++++++++ V
%: -Code.
\hCodeTag{terminal}
\begin{lstlisting}[style=terminalstyle, frame=single, rulecolor=\color{termframe}, backgroundcolor=\color{termbg}, basicstyle=\ttfamily\footnotesize\color{termfg}]
$ java -jar MyApp.jar aaa
~~MyClass
i=0:aaa
\end{lstlisting}
% +++++++++++++++++++++++++++++++++++++++ A


% ~~~~~~~~~~~~~~~~~~~~~~~~~~~~~~~~~~~~~~
\subsection{Inline code}

% +++++++++++++++++++++++++++++++++++++++ V
%: -Code.
\begin{texcode}
\hCode{MyClass.java}             % plain, colored, 2 pt smaller: escape _ { } \ as in text
\hCode{project\_dir}             % gives project_dir
\hCodeJ{int n = s.length();}     % Java, syntax colored, typed as is
\hCodeB{echo 'aa bbb' >file}     % bash, typed as is (quotes stay straight)
\end{texcode}
% +++++++++++++++++++++++++++++++++++++++ A

Results:
\hCode{MyClass.java},
\hCode{project\_dir},
\hCodeJ{int n = s.length();},
\hCodeB{echo 'aa bbb' >file}.

\hNote{\hCode{\textbackslash hCodeJ} and \hCode{\textbackslash hCodeB} cannot be used in section titles,
captions or footnotes, and their text must not contain \hCode{\%} or \hCode{\#}.
Use \hCode{\textbackslash hCode} there, with the special characters escaped.}




% ~~~~~~~~~~~~~~~~~~~~~~~~~~~~~~~~~~~~~~
\section{Lists, tables, trees, figures}

% ~~~~~~~~~~~~~~~~~~~~~~~~~~~~~~~~~~~~~~
\subsection{Lists (enumitem)}

% +++++++++++++++++++++++++++++++++++++++ V
%: -Code.
\begin{texcode}
\begin{enumerate}[label=\roman*., itemsep=0pt]   % i., ii.; no extra space between items
	\item first
	\item second
\end{enumerate}
\end{texcode}
% +++++++++++++++++++++++++++++++++++++++ A

\begin{enumerate}[label=\roman*., itemsep=0pt]
	\item first
	\item second
\end{enumerate}

Other options: \hCode{[nosep]} (no space at all), \hCode{[label=(\textbackslash alph*)]} for (a), (b),
\hCode{[topsep=1pt]}.


% ~~~~~~~~~~~~~~~~~~~~~~~~~~~~~~~~~~~~~~
\subsection{Menus and keys (menukeys)}

% +++++++++++++++++++++++++++++++++++++++ V
%: -Code.
\begin{texcode}
\menu{File > Save As}    \keys{Ctrl + C}    \keys{\shift + \cmd + B}
\end{texcode}
% +++++++++++++++++++++++++++++++++++++++ A

\menu{File > Save As} \quad \keys{Ctrl + C} \quad \keys{\shift + \cmd + B}


% ~~~~~~~~~~~~~~~~~~~~~~~~~~~~~~~~~~~~~~
\subsection{Folder trees (forest)}

% +++++++++++++++++++++++++++++++++++++++ V
%: -Code.
\begin{texcode}
\begin{forest}
for tree={font=\ttfamily\small, grow'=0, folder, draw=none,
          s sep=0pt, inner ysep=1pt}            % less space between lines
	[project/
		[src/
			[packageA/ [modifyMe.java] [doNotModifyMe.java]]
		]
	]
\end{forest}
\end{texcode}
% +++++++++++++++++++++++++++++++++++++++ A

\begin{center}
\begin{forest}
for tree={font=\ttfamily\small, grow'=0, folder, draw=none, s sep=0pt, inner ysep=1pt}
	[project/
		[src/
			[packageA/ [modifyMe.java] [doNotModifyMe.java]]
		]
	]
\end{forest}
\end{center}

\hNote{\hCode{s sep} is the space between lines, \hCode{inner ysep} the padding around each name.
Do not change \hCode{l sep}: with \hCode{folder} it breaks the connecting lines.
\hCode{folder indent=.5em} makes the indent smaller.}


% ~~~~~~~~~~~~~~~~~~~~~~~~~~~~~~~~~~~~~~
\subsection{Algorithms (algorithm2e)}

% +++++++++++++++++++++++++++++++++++++++ V
%: -Code.
\begin{texcode}
\begin{algorithm}[H]
	\caption{isBalanced($s$)}\label{alg:balanced}
	\KwIn{a string $s$}
	\KwOut{true or false}
	$S \gets$ an empty stack\;
	\ForEach{character $c$ of $s$}{
		\lIf{$c$ is an opening bracket}{push $c$ on $S$}
		\ElseIf{$c$ is a closing bracket}{
			\lIf{$S$ is empty \textbf{or} pop($S$) does not match $c$}{\Return{false}}
		}
	}
	\Return{$S$ is empty}\;
\end{algorithm}
\end{texcode}
% +++++++++++++++++++++++++++++++++++++++ A

\begin{algorithm}[H]
	\caption{isBalanced($s$)}\label{alg:balanced}
	\KwIn{a string $s$}
	\KwOut{true or false}
	$S \gets$ an empty stack\;
	\ForEach{character $c$ of $s$}{
		\lIf{$c$ is an opening bracket}{push $c$ on $S$}
		\ElseIf{$c$ is a closing bracket}{
			\lIf{$S$ is empty \textbf{or} pop($S$) does not match $c$}{\Return{false}}
		}
	}
	\Return{$S$ is empty}\;
\end{algorithm}

\refalg{alg:balanced} is placed here because of \hCode{[H]}.
Without it, it floats like a figure.

\hNote{\hCode{\textbackslash clearpage} prints all waiting figures, tables and algorithms, then starts a new page.
To print them without a new page, use \hCode{\textbackslash FloatBarrier} of the \hCode{placeins} package.}




% ~~~~~~~~~~~~~~~~~~~~~~~~~~~~~~~~~~~~~~
\section{Section numbering}\label{sec:numbering}

Sections are numbered 1, 2, \dots\ with subsections 1.1 and subsubsections 1.1.1.
The style can switch before any \hCode{\textbackslash section}:

% +++++++++++++++++++++++++++++++++++++++ V
%: -Code.
\begin{texcode}
\hNumbersLetter{G}        % next section: G, subsections G1, G2, subsubsections G1.1
\section{Group work}
\hNumbersArabic           % back to numbers; they continue where they stopped
\hNumbersArabic[1]        % back to numbers, restarting at 1
\end{texcode}
% +++++++++++++++++++++++++++++++++++++++ A

Lab handouts use four ready-made sections:

% +++++++++++++++++++++++++++++++++++++++ V
%: -Code.
\begin{texcode}
\hSectionWarmUp      % W  WarmUp
\hSectionGoal        % G  Goal
\hSectionChallenge   % C  Challenge
\hSectionHome        % H  Home
\end{texcode}
% +++++++++++++++++++++++++++++++++++++++ A

The next four sections are made with them, and the section after them with \hCode{\textbackslash hNumbersArabic}.
Look at their numbers here and in the contents.

\hNote{\hCode{\textbackslash ref} prints the number in the current style, e.g.\ \refsec{sec:goalA} and \refsec{sec:numbering}.
Links in the PDF stay correct, even though G and a later section 7 have the same counter value.}



%: ~~~~~~~~~~~~~~~~~~~~~~~~~~~~~~~~~~~~~~ V WarmUp
\hSectionWarmUp
\subsection{A warm-up task}
Complete \hCode{lowerFirstChar}.
% ~~~~~~~~~~~~~~~~~~~~~~~~~~~~~~~~~~~~~~ A WarmUp



%: ~~~~~~~~~~~~~~~~~~~~~~~~~~~~~~~~~~~~~~ V Goal
\hSectionGoal
\subsection{First goal}\label{sec:goalA}
\subsubsection{A detail}
\subsection{Second goal}
% ~~~~~~~~~~~~~~~~~~~~~~~~~~~~~~~~~~~~~~ A Goal



%: ~~~~~~~~~~~~~~~~~~~~~~~~~~~~~~~~~~~~~~ V Challenge
\hSectionChallenge
\subsection{A challenge}
% ~~~~~~~~~~~~~~~~~~~~~~~~~~~~~~~~~~~~~~ A Challenge



%: ~~~~~~~~~~~~~~~~~~~~~~~~~~~~~~~~~~~~~~ V Home
\hSectionHome
\subsection{At home}
% ~~~~~~~~~~~~~~~~~~~~~~~~~~~~~~~~~~~~~~ A Home




% ~~~~~~~~~~~~~~~~~~~~~~~~~~~~~~~~~~~~~~
\hNumbersArabic
\section{Macros}

% ~~~~~~~~~~~~~~~~~~~~~~~~~~~~~~~~~~~~~~
\subsection{References}

% +++++++++++++++++++++++++++++++++++++++ V
%: -Code.
\begin{texcode}
\reffig{fig:x}  \refeq{eq:x}  \reftbl{tbl:x}  \refsec{sec:x}  \refalg{alg:x}  \reflst{lst:x}
\refthm{thm:x}  \refthmA{thm:x}{item:b}  \reflem{lem:x}  \refdef{def:x}  \refexmp{ex:x}
\refcite{key}
\end{texcode}
% +++++++++++++++++++++++++++++++++++++++ A

They print ``Fig.~3'', ``Eq.~2'', ``Table~1'', ``Sec.~4'', ``Alg.~1'', ``List~2'',
``Theorem~1'', ``Theorem~1(b)'', ``Lemma~2'', ``Definition~3'', ``Example~4'', ``Ref.~[7]'',
with a non-breaking space. In this document:
\reftbl{tbl:styles}, \refalg{alg:balanced}, \refsec{sec:tags}.


% ~~~~~~~~~~~~~~~~~~~~~~~~~~~~~~~~~~~~~~
\subsection{Text}

% +++++++++++++++++++++++++++++++++++++++ V
%: -Code.
\begin{texcode}
\hQuote{Talk is cheap. Show me the code.}
\hIdea{add an example here}
text\hFootnote{HB}{check this}
\hAbs{x - y}
\YULearn
\hFontLarger{2}{two points larger}  \hFontLarger{-2}{two points smaller}
\end{texcode}
% +++++++++++++++++++++++++++++++++++++++ A

\hQuote{Talk is cheap. Show me the code.}
\quad \hIdea{add an example here}
\quad text\hFootnote{HB}{check this}
\quad \hAbs{x - y}
\quad \YULearn
\quad \hFontLarger{2}{two points larger}, \hFontLarger{-2}{two points smaller}

\hNote{A document may define its own \hCode{\textbackslash YULearn} with \hCode{\textbackslash newcommand};
the library then uses the document's.}


% ~~~~~~~~~~~~~~~~~~~~~~~~~~~~~~~~~~~~~~
\subsection{Time stamps}

% +++++++++++++++++++++++++++++++++++++++ V
%: -Code.
\begin{texcode}
\hbTimeStamp        % in the footer
\hbTimeStampDraft   % red, with the time: for drafts
\end{texcode}
% +++++++++++++++++++++++++++++++++++++++ A

\hbTimeStamp \quad \hbTimeStampDraft




% ~~~~~~~~~~~~~~~~~~~~~~~~~~~~~~~~~~~~~~
\section{Layout}

% ~~~~~~~~~~~~~~~~~~~~~~~~~~~~~~~~~~~~~~
\subsection{Titles and headings}

With \hCode{amsart}, the title, the authors and the section headings are printed as typed,
not in capitals or small caps.
\hCode{\textbackslash usepackage[amsheadings]\{hblibArticle\}} keeps the \hCode{amsart} look.


% ~~~~~~~~~~~~~~~~~~~~~~~~~~~~~~~~~~~~~~
\subsection{Table of contents}

One size smaller than the text; sections in bold; subsections indented, subsubsections more.

% +++++++++++++++++++++++++++++++++++++++ V
%: -Code.
\begin{texcode}
\setcounter{tocdepth}{2}   % after \usepackage{hblibArticle}: no subsubsections in the contents
\end{texcode}
% +++++++++++++++++++++++++++++++++++++++ A


% ~~~~~~~~~~~~~~~~~~~~~~~~~~~~~~~~~~~~~~
\subsection{Footer}

Every page, also the first one with \hCode{amsart}:
a rule, then on the left the file name and the date, on the right ``page / pages''.
The file name is \hCode{\textbackslash jobname}, the name of the .tex file.


% ~~~~~~~~~~~~~~~~~~~~~~~~~~~~~~~~~~~~~~
\subsection{Pages and columns}

% +++++++++++++++++++++++++++++++++++++++ V
%: -Code.
\begin{texcode}
\newpage     % a new page (in two columns: a new column)
\clearpage   % a new page, after printing all waiting figures, tables and algorithms
\documentclass[11pt, twocolumn]{article}   % two columns
\end{texcode}
% +++++++++++++++++++++++++++++++++++++++ A


% ~~~~~~~~~~~~~~~~~~~~~~~~~~~~~~~~~~~~~~
\subsection{Hyperlinks}

The contents, \hCode{\textbackslash ref} and URLs are clickable.
A red box marks an internal link, a blue box a URL; the boxes are shown on screen, not printed.

\url{https://github.com/halukbingol/LaTeX-Templates}

% +++++++++++++++++++++++++++++++++++++++ V
%: -Code.
\begin{texcode}
\hypersetup{hidelinks}                       % after \usepackage{hblibArticle}: no marks
\hypersetup{colorlinks=true, linkcolor=blue!70!black, urlcolor=blue, citecolor=green!50!black}
                                             % colored text instead of boxes
\end{texcode}
% +++++++++++++++++++++++++++++++++++++++ A

\hNote{\hCode{\textbackslash hCode} in a section title prints correctly, but hyperref warns
``Token not allowed in a PDF string'' for the PDF bookmark. The warning is harmless.}




% ~~~~~~~~~~~~~~~~~~~~~~~~~~~~~~~~~~~~~~
\section{Customizing}\label{sec:customizing}

Put these after \cs{usepackage}\marg{hblibArticle}; inside a group \hCode{\{...\}} they change only that part.

% +++++++++++++++++++++++++++++++++++++++ V
%: -Code.
\begin{texcode}
\geometry{margin=2.5cm}                    % margins
\definecolor{colorHCode}{RGB}{0,90,200}    % \hCode in blue
\definecolor{codekw}{RGB}{160,0,80}        % keywords in purple
\setcounter{tocdepth}{2}                   % no subsubsections in the contents
\hCodeTagsfalse                            % no language tags
\renewcommand{\hbTimeStamp}{v1.2}          % the version in the footer
\end{texcode}
% +++++++++++++++++++++++++++++++++++++++ A

One listing only: give the option in brackets.

% +++++++++++++++++++++++++++++++++++++++ V
%: -Code.
\begin{texcode}
\begin{javacode}[basicstyle=\ttfamily\scriptsize, numbers=left]
int x = 1;   // smaller, numbered
\end{javacode}
\end{texcode}
% +++++++++++++++++++++++++++++++++++++++ A

% +++++++++++++++++++++++++++++++++++++++ V
%: -Code.
\begin{javacode}[basicstyle=\ttfamily\scriptsize, numbers=left]
int x = 1;   // smaller, numbered
\end{javacode}
% +++++++++++++++++++++++++++++++++++++++ A

A color only in a group:
{\definecolor{colorHCode}{RGB}{0,90,200}\hCode{blue here}}, \hCode{the usual color here}.




% ~~~~~~~~~~~~~~~~~~~~~~~~~~~~~~~~~~~~~~
\section{Troubleshooting}\label{sec:troubleshooting}

\begin{center}\small
\begin{tabular}{@{}L{0.34\linewidth}L{0.28\linewidth}L{0.32\linewidth}@{}}
	\toprule
	\textbf{Message or problem} & \textbf{Cause} & \textbf{Fix} \\
	\midrule
	Missing \textbackslash begin\{document\} right after loading
		& \cs{input}\marg{hblibCommon.sty} & \cs{usepackage}\marg{hblibCommon} \\
	File `hblibCommon.tex' not found
		& \cs{input}\marg{hblibCommon} & \cs{usepackage}\marg{hblibCommon} \\
	You have requested package `../../\dots/hblibCommon', but the package provides `hblibCommon'
		& a relative path in \cs{usepackage} & harmless \\
	Footer shows ``1 / ??''; contents or references empty or ``??''
		& first run & compile again \\
	Command \textbackslash hNumbersArabic already defined
		& the document still has its own copy of the numbering commands & delete the copy \\
	Command \textbackslash code (or another name) already defined
		& the document defines a name the library has & rename it, or use \cs{renewcommand} \\
	Token not allowed in a PDF string
		& \cs{hCode} in a section title & harmless \\
	Package Listings Error: style codesnippet undefined
		& an old library (hblibSimple) & use hblibCommon \\
	\cs{hCodeJ} fails in a title, caption or footnote
		& listings cannot work inside another command's argument & \cs{hCode} there \\
	\cs{javafile} with \hCode{firstline=5} is numbered 1, 2, 3
		& listings counts the printed lines & add \hCode{firstnumber=5} \\
	the tag covers the end of the first code line
		& a full-width first line & shorten it, or \cs{hCodeTagsfalse} before the listing \\
	TeX capacity exceeded, sorry [input stack size=10000] at \cs{maketitle} (amsart)
		& a size command such as \cs{normalsize} in \cs{title}; amsart copies the title to the page header & give a plain short title: \cs{title}\oarg{Short}\marg{Long \textbackslash\textbackslash\ \{\textbackslash small \dots\}} \\
	a figure or algorithm appears pages later
		& floats wait for room & \oarg{H}, \cs{clearpage}, or \cs{FloatBarrier} (placeins) \\
	\bottomrule
\end{tabular}
\end{center}



\end{document}
